% !TeX program = xelatex
\documentclass{UAmathtalk}
\usepackage{setspace}
\renewcommand{\when}{12:00 p.m.}
%\renewcommand{\where}{Hudson 0110. Zoom}
\renewcommand{\where}{Hudson 0110.\\ \href{https://albany.zoom.us/j/81219478289?pwd=N5fy7vxKBXqlQ4mAq4KaeaAM5iAHi4.1}{Zoom ID: 812\,1947\,8289 - Passcode: 233444}}
\author{Roy Nicolas Nehme}
\address{LAGA - Université Paris 13}
%\urladdr{https://www.smajhi.com/}
\title{Pruning Distance of Upset-Decomposable Persistence Modules}
\date{Friday, September 25, 2026}


\begin{document}

\maketitle
%\renewcommand*{\abstractsize}{\normalsize}
\renewcommand*{\abstractsize}{\fontsize{13.5}{16.2}\selectfont}
\begin{abstract}
{Multiparameter persistence offers a framework for studying data that depends on several parameters. However, multiparameter persistence modules are not described in a stable way by their indecomposable decompositions, and classical distances either unstable or difficult to compute. This motivates the study of alternative distances that are both computable and stable.

In this talk, we study the pruning distance, recently introduced by Håvard Bjerkevik, for upset-decomposable persistence modules. We prove that, for modules of finite pointwise dimension, the pruning distance is stable, in the sense that it is Lipschitz equivalent to the interleaving distance.

Our proof is based on an explicit computation of the pruning of an upset-decomposable module, followed by a combinatorial argument using directed graphs.}
\end{abstract}
\end{document}


